\section{The arithmetic mass count}\label{sec:arithmetic}
This section evaluates the count \eqref{eq:Darith}. We write each
$d>1$ as its largest prime factor times a squarefree number with smaller
prime factors, count the latter with smooth-number estimates, and sum over
the largest prime. Because the classical estimates are uniform only in
limited ranges, the sum is split into three ranges of the largest prime.

For a prime $p$, put
\[
 m(p)=\prod_{q\le p}\frac q{q+1}.
\]
Let $\Psi(x,y)$ count the positive integers $n\le x$ with $P^+(n)\le y$.
Use $\Psi(x,p^-)$ for the strict condition $P^+(n)<p$, and let
$\Psi_{\mathrm{sf}}(x;p^-)$ impose squarefreeness as well.
All three functions are zero for $x<1$.

\begin{theorem}\label{thm:arithmetic}
For the masses $\nu_d$ of Section~\ref{sec:intro},
\[
 D(t)=\#\{d\in\sfq:\nu_d>t\}
 \sim\frac1{t\log(1/t)}\qquad(t\downarrow0).
\]
\end{theorem}

\subsection{Classical inputs and the uniformity actually used}
We use Chebyshev's upper bound $\pi(y)\ll y/\log y$ and the two Mertens
asymptotics
\[
 \prod_{p\le y}(1-1/p)\sim\frac{e^{-\gamma}}{\log y},
 \qquad \sum_{p\le y}\frac1p=\log\log y+B+o(1).
\]
In particular,
\begin{equation}\label{eq:mproduct}
 m(p)\sim\frac{\zeta(2)e^{-\gamma}}{\log p},
 \qquad \frac{c}{\log p}\le m(p)\le\frac C{\log p}\quad(p\ge2)
\end{equation}
for absolute positive constants $c,C$.

The Dickman function $\rho$ is $1$ on $[0,1]$, zero on $(-\infty,0)$, and
satisfies $u\rho'(u)=-\rho(u-1)$ for $u>1$. We use:
\begin{align}
 \Psi(x,y)&\ll x e^{-u/2}
 &&(x\ge y\ge2,\ u=\log x/\log y),\label{eq:rankin}\\
 \sup_{u\in I}\left|\frac{\Psi(x,x^{1/u})}{x}-\rho(u)\right|&\longrightarrow0
 &&(x\to\infty),\label{eq:dickmancompact}\\
 \int_0^\infty\rho(u)\,du&=e^\gamma.\label{eq:dickmanintegral}
\end{align}
In \eqref{eq:dickmancompact}, $I$ is any fixed compact interval in
$(0,\infty)$.

These facts may be sourced to Hildebrand--Tenenbaum~\cite{HT}.
Printed page~414 states the fixed-$u$ Dickman limit. Equation~(1.4) there
has an extra factor $\log y$ in \eqref{eq:rankin}; the following sentence
explicitly states that this factor can be removed, citing Tenenbaum's
Theorem~III.5.1. Thus \eqref{eq:rankin} uses that strengthened bound,
not equation~(1.4) alone. For \eqref{eq:dickmanintegral}, use Lemma~2.6
on printed page~426 and the definition of $E$ in (2.3), printed page~422:
the Laplace transform is $\exp(\gamma+E(-s))$, and $E(0)=0$.

For completeness, the compact uniformity in \eqref{eq:dickmancompact}
follows from the fixed-$u$ limit. Both
$u\mapsto\Psi(x,x^{1/u})/x$ and $\rho(u)$ are nonincreasing; $\rho$ is
continuous. A finite mesh on $I$ with small oscillation of $\rho$ brackets
all intermediate values between endpoint values, where convergence is
simultaneous. For $u<1$ the smooth count is simply $\lfloor x\rfloor$.
No approximation on a growing $u$-range is used below.

\subsection{An exact weak-count identity}
Let
\[
 F(X)=\#\{d\in\sfq:\nu_d\ge1/X\}\qquad(X>1).
\]
Every $d>1$ factors uniquely as $d=pm$ with $p=P^+(d)$ and
$m$ squarefree with $P^+(m)<p$. Thus
\begin{equation}\label{eq:Fidentity}
 F(X)=1+\sum_p\Psi_{\mathrm{sf}}(x_p;p^-),
 \qquad x_p=\frac{Xm(p)}p.
\end{equation}
Primes $p\ge X$ give zero. We prove $F(X)\sim X/\log X$, then transfer to
the strict count. Write $L=\log X$ and $\alpha_p=(\log p)/L$.
Fix $0<a\le1/4<b<1$ and split the sum at $X^a$ and $X^b$.

\subsection{Small largest primes}
The term $p=2$ equals $1$ for all sufficiently large $X$, since $m=1$
is the only squarefree positive integer with all prime factors below $2$.
For $3\le p\le X^a$, \eqref{eq:mproduct} implies, uniformly in this range,
\[
 x_p\ge \frac{cX^{1-a}}{L}\ge p,\qquad \log x_p\ge L/2
\]
once $X$ is large. Hence \eqref{eq:rankin} gives
\[
 \Psi_{\mathrm{sf}}(x_p;p^-)
 \le\Psi(x_p,p)
 \ll\frac{X}{p\log p}\exp\left(-\frac{L}{4\log p}\right).
\]
For integer $k\ge1$, Chebyshev's bound gives
\[
 \sum_{e^k<p\le e^{k+1}}\frac1p
 \le e^{-k}\pi(e^{k+1})\ll\frac1k.
\]
On such a bin, $1/\log p<1/k$ and
$e^{-L/(4\log p)}\le e^{-L/(4(k+1))}\le e^{-L/(8k)}$, since $4(k+1)\le8k$
for $k\ge1$. The small-prime contribution, apart from $p=2$, is therefore
at most
\[
 CX\sum_{1\le k\le aL}\frac1{k^2}e^{-L/(8k)}.
\]
The function $f(u)=u^{-2}e^{-L/(8u)}$ is unimodal, with maximum
$256e^{-2}/L^2$. Comparing the sum with its integral and twice its maximum,
and using
\[
 \int_0^{aL}u^{-2}e^{-L/(8u)}\,du
 =\frac8L e^{-1/(8a)},
\]
we obtain
\begin{equation}\label{eq:smallrange}
 \sum_{p\le X^a}\Psi_{\mathrm{sf}}(x_p;p^-)
 \ll \frac XL e^{-1/(8a)}+\frac X{L^2}+1.
\end{equation}
The implied constant here can be taken independent of $a$ in $(0,1/4]$.

\subsection{The middle range}
We first establish, uniformly for $X^a<p<X^b$,
\begin{equation}\label{eq:squarefree-smooth}
 \Psi_{\mathrm{sf}}(x_p;p^-)
 =x_p\{c_0\rho(u_p)+o(1)\},
 \qquad u_p=\frac{\log x_p}{\log p}.
\end{equation}
The squarefree identity $\mu^2(n)=\sum_{\beta^2\mid n}\mu(\beta)$ gives
\[
 \Psi_{\mathrm{sf}}(x;p^-)
 =\sum_{\substack{\beta\le\sqrt x\\P^+(\beta)<p}}
   \mu(\beta)\Psi(x/\beta^2,p^-).
\]
Terms with $\beta>B$ have absolute sum at most
$x\sum_{\beta>B}\beta^{-2}\ll x/B$. For fixed $B$, all $\beta\le B$
satisfy $P^+(\beta)<p$ once $X$ is large. Uniformly in the middle range,
$x_p/\beta^2\to\infty$ and
\[
 \frac{\log(x_p/\beta^2)}{\log p}
\]
stays in a fixed compact subinterval of $(0,\infty)$.
Replacing the strict cutoff $p^-$ by $p$ changes a count at argument $z$
by at most $z/p$. Apply \eqref{eq:dickmancompact}, and then continuity
of $\rho$, to obtain
\[
 \frac{\Psi(x_p/\beta^2,p^-)}{x_p}
 =\beta^{-2}\rho(u_p)+o(1)
\]
uniformly for this finite set of $\beta$. First let $X\to\infty$ and then
$B\to\infty$. Since $\sum_{\beta\ge1}\mu(\beta)/\beta^2=1/\zeta(2)=c_0$,
this proves \eqref{eq:squarefree-smooth}.

Equation~\eqref{eq:mproduct} now gives, uniformly on the middle range,
\[
 u_p=\frac{1-\alpha_p}{\alpha_p}+o(1),\qquad
 c_0x_p=\frac{e^{-\gamma}X}{p\log p}\{1+o(1)\}.
\]
The total mass $\sum_{X^a<p<X^b}x_p$ is $O_{a,b}(X/L)$, so the uniform
errors remain $o_{a,b}(X/L)$ after summation. Consequently
\begin{equation}\label{eq:middlebefore}
 \frac LX\sum_{X^a<p<X^b}\Psi_{\mathrm{sf}}(x_p;p^-)
 =e^{-\gamma}\sum_{X^a<p<X^b}
 \frac1p\,\frac{\rho((1-\alpha_p)/\alpha_p)}{\alpha_p}+o(1).
\end{equation}
Mertens' reciprocal-prime asymptotic implies weak convergence of these
prime measures on $[a,b]$ to $d\alpha/\alpha$: on any fixed subinterval
$[u,v]\subset[a,b]$ their masses tend to $\log(v/u)$.
Approximating the continuous integrand in \eqref{eq:middlebefore} by
step functions yields
\begin{equation}\label{eq:middlelimit}
 \lim_{X\to\infty}\frac LX
 \sum_{X^a<p<X^b}\Psi_{\mathrm{sf}}(x_p;p^-)
 =e^{-\gamma}\int_a^b
     \rho\!\left(\frac{1-\alpha}{\alpha}\right)\frac{d\alpha}{\alpha^2}
 =e^{-\gamma}\int_{(1-b)/b}^{(1-a)/a}\rho(v)\,dv.
\end{equation}

\subsection{Large largest primes and completion}
The trivial estimate $\Psi_{\mathrm{sf}}(x_p;p^-)\le x_p$ and
\eqref{eq:mproduct} show
\begin{equation}\label{eq:largerange}
 \sum_{X^b\le p<X}\Psi_{\mathrm{sf}}(x_p;p^-)
 \le \frac{CX}{bL}\sum_{X^b\le p<X}\frac1p
 =\frac{CX}{bL}\{\log(1/b)+o(1)\}.
\end{equation}
All summands in \eqref{eq:Fidentity} are nonnegative. Thus the limiting
middle contribution is a lower bound for $\liminf LF(X)/X$, while
\eqref{eq:smallrange} and \eqref{eq:largerange} bound the excess in the
upper limit. Let $a\downarrow0$ and $b\uparrow1$, after taking the limits
in $X$. Equation~\eqref{eq:dickmanintegral} gives
\[
 \frac{LF(X)}X\longrightarrow
 e^{-\gamma}\int_0^\infty\rho(v)\,dv=1.
\]
Finally, for every fixed $\varepsilon>0$,
\[
 F\bigl(X/(1+\varepsilon)\bigr)\le D(1/X)\le F(X).
\]
Taking $X\to\infty$ and then $\varepsilon\downarrow0$ proves
Theorem~\ref{thm:arithmetic}. In particular, $D$ is regularly varying
at zero with index $-1$.
